% !TeX root = ../papers/nc-p3.tex
\section{四次Grassmann写像と空間楕円八次曲線による復元}
\label{sec:finite-quartic-grassmann-windows}
\begin{theorem}\label{thm:finite-quartic-grassmann-windows}
反対階数10の有限片側対称窓で、二つの奇位置から得る層を $E,E'$ とする。
両方が局所自由なら $E(2)$ はちょうど四つの独立な大域切断で生成され、
その完全評価列は
\[
 0\to E'^*(-2)\to H\otimes\mathcal O\to E(2)\to0,\qquad \dim H=4
\]
である。従って全ての正則な対象は、二平面核の写像
$f:\mathbb P(U)\to\mathrm{Gr}(2,H)$、$f^*\mathcal O(1)=\mathcal O(4)$ から復元できる。
二窓の添字保持同型はこの写像の始域の射影変更と $H$ の線型変更による同値と必要十分である。
二移動は $f(x)\mapsto f(x)^\perp\subset H^*$ に対応し、正則窓の最小周期は
この双対写像と同値なら2、そうでなければ4である（対称軸の位相を固定する）。
\end{theorem}
\begin{proof}
定理\ref{thm:finite-symmetric-instanton-monad}の極小分解を層化して1捩りする。
最後の二項の余核は、反対の三段加群が次数を反転した双対であることより
$(\ker b')^*(-4)$ と同定できる。中央項への包含から $4\mathcal O(-3)$ への射影は、
反対の線型部分 $a'$ を双対化し $-4$ 捩りした写像である。$a'$ が部分束なので全射で、核は $E'^*(-4)$ である。
中央項の残り $4\mathcal O(-2)\oplus4\mathcal O(-1)$ から $\ker b$ への写像は従って全射で、
核が今の $E'^*(-4)$ となる。さらに $4\mathcal O(-1)$ の写像は $a$ であり、
これを割って2捩りすると表示の列を得る。$H^0(E'^*(-2))=H^1(E'^*(-2))=0$、
かつ行列式は自明なので、これは大域切断全体の評価である。核の双対の評価は他方の完全評価列となる。
$f$ から普遍商の引戻しと2捩り戻しで $E$ が得られ、元の窓を復元する。
$\det E(2)=\mathcal O(4)$ は豊富なので $f$ はその像への有限射である。
二移動は二層を交換する。正則な両対称窓は定理\ref{thm:all-bisymmetric-as-closure}で
反対階数8以下なので、奇数周期で他方へ対称軸が移ることはない。
\end{proof}

\begin{theorem}\label{thm:finite-elliptic-octic-windows}
同じ両局所自由部分は、滑らかな楕円曲線 $C$、次数8の線束 $L$、
$\mathbb P^3$ への閉埋込みを与える $D\subset H^0(L)$（次元4）からの次の開条件付き構成で尽くされる。
\[
 N_j= H^0(L^{j+1})/D^{j+1}\quad (0\le j\le2)
\]
を $D$ の乗法で結び、次元 $4,6,4$、二回乗法に対する有限Frobenius検査・有限辺・反対階数10、
両方のモナドの部分束条件を課す。さらに四次の評価の全射性と、次数1の関係による表示が $N_2$ を与える条件を課す。ここで $D^m$ は積の像である。
同じ窓の楕円八次曲線は射影変更を除いて $E(2)$ の滑らかな零曲線を持つ切断の
$\mathbb P^3$ の開部分で尽きる（束を保つ射影座標変更も作用する）。
この有限候補部分の添字保持同型類は既約で次元14であり、その幾何学的一般点は全次数で正則である。
特にこの位相の反対階数10の有限片側対称\emph{正則窓全体の閉包}は既約である。
\end{theorem}
\begin{proof}
大域生成された $E(2)$ の一般切断の零点 $C$ は滑らかで余次元2である。
実際 $c_2(E(2))=8$ で空ではなく、普遍零点は評価核の束なので標数零の一般滑らか性を使える。
随伴公式は $\omega_C=\mathcal O_C$ を与える。また
\[
 0\to\mathcal O\to E(2)\to I_C(4)\to0
\]
を $-4$ 捩りした列から $H^1(I_C)=H^1(E(-2))=0$、従って $C$ は連結な楕円八次曲線である。
同じ列と直線束のコホモロジーから全次数で $H^1(E(j-1))=H^1(I_C(j+1))$ を得る。
後者は記載の切断の商そのもので、乗法と両立する。

逆に曲線データで三つの次元条件を満たせば $H^0(I_C(1)),H^0(I_C(2)),H^0(I_C(3))$ は零である。
さらに曲線のSerre構成で元の切断を復元できることを直接確認する。
$\operatorname{Ext}^1(I_C(4),\mathcal O)\simeq
H^0(\mathcal Ext^2(\mathcal O_C(4),\mathcal O))=H^0(\omega_C)=k$
である。実際 $H^1,H^2$ の中間直線束コホモロジーは零で、局所完全交叉の他の局所Extは零である。
非零類は各点で二方程式のKoszul表示の単元類なので、拡大の中項は局所自由である。
その $-2$ 捩りを $E_C$ とすれば、行列式は自明で $H^0(E_C)=0$ なので安定である。
$H^1(E_C(-2))=H^1(I_C)=0$。ただし三つの隣合う次数だけから全Rao加群との一致を
仮定しないため、逆構成の開条件には $H^1(I_C(4))=0$ と
$D\otimes N_0\to N_1$ の核による一次表示が $N_2$ を与えることも含める。
順方向ではこれらは既に成立する。
$I_C$ は5正則である（$H^2(I_C(3))=H^1(L^3)=0, H^3(I_C(2))=0$）。
これで高次のRaoは零、負次数と次数0も連結性と負次数切断の不在で零となる。
したがって $H^1_*(E_C)$ は復元した三段加群そのものである。

この加群の線型モナドが与える束 $E$ と $E_C$ の一致は次の初等的な構成でも分かる。
$H^1(E_C(-1))$ の基底で普遍拡大 $0\to E_C\to J\to4\mathcal O(1)\to0$ を取ると、
全次数で乗法がこの四元で生成するので $H^1_*(J)=0$。
双対側でも $E_C^*=E_C$ であるから、$H^1_*(J^*)=H^1_*(E_C)$ を消す四元の
普遍拡大を双対化して $0\to4\mathcal O(-1)\to B\to J\to0$ を得る。
$H^1_*(B)=H^2_*(B)=0$ なのでHorrocksの分裂判定で $B$ は直線束の直和である。
$B\to4\mathcal O(1)$ の切断上の像は三段加群の関係を生成する。
その極小生成は次数1に十本で、$B$ の階数も10なので、十本全てが $\mathcal O$ から来る。
複体の左端も四つの線型シジジーを張るので $E_C=E$ である。
ここで $J^*\to E_C^*$ の第一コホモロジーの同型には $4\mathcal O(-1)$ の中間消滅を用いた。
逆に同じ窓ではこの束が決まり、拡大の一意性により曲線の選択は表示の切断だけである。

基礎の楕円曲線と次数8の線束の同型類は一次元既約である（平行移動で任意の二つの次数8の束が移る）。
その上の $\mathrm{Gr}(4,8)$ の開部分なので、曲線表示の同型類は既約で次元17である。
原点を選んで $L=\mathcal O(8o)$ としても残余自己同型は有限で、同じ次元となる。
開条件は空でない。定理\ref{thm:twenty-point-distinct-companion}の例の両辺は
$B_-=-B_+(J\otimes J)$ で移り、前節のモナドの部分束検算が両側で通る。
その窓の二十点は被約で、その位置を全て固定する射影変換は自明なので一般安定化群は有限である。
よって一般ファイバーの次元3を引いて14となる。
標準複体の各固定内部次数の完全性は開であり、この正則例で全て成立するため
既約なパラメータの幾何学的一般点では各次数で完全である。個別の閉点での正則性はここでは従わない。
\end{proof}
